\chap{Abstract} 


%%%% your abstract goes here (word limit: 350)
% DRAFT (2026-07-28) -- edit to final voice; ~330 words, under the 350 limit.
% Numbers trace to Chapter 3 / results_hpc. No citations in abstract (per convention).
Premalignant epithelial lesions progress within tissue ecosystems whose inflammatory and stromal
composition changes across the precursor sequence, yet whether the local microenvironment alters an
epithelial cell's progression-aligned regulatory state beyond its own intrinsic state---and at what
spatial scale---has been difficult to isolate, because cross-sectional profiling captures snapshots
rather than trajectories and confounds cell state with tissue context. This dissertation develops
StageBridge, a biologically structured universal differential equation that separates an intrinsic
regulatory-progression field from a niche-gated context field and defines a per-program, conditional
context residual, $R_{\text{cond}}$, as the additional regulatory displacement attributable to local
ecology within ecologically comparable cells. The transport coupling underlying the estimand was
selected against a frozen synthetic identifiability benchmark, which falsified an initial
state-neutral coupling and motivated the conditional-stratified formulation used throughout. Applied to a patient-held-out
lung adenocarcinoma precursor cohort spanning atypical adenomatous hyperplasia, adenocarcinoma in
situ, and invasive disease, StageBridge recovered a small, reproducible set of context-associated
regulatory programs---twenty-nine on the preinvasive AAH-to-AIS edge and twenty-eight at the
AIS-to-invasive boundary---forming a coordinated WNT, p53, TRAIL, and AP-1 axis with coupled interferon
and TGF$\beta$-effector suppression, sign-consistent across all folds and donors and concordant with
co-expression analysis. A pre-registered falsification ladder then showed that this
contextual information did not resolve at the single-cell-neighborhood scale: cell-resolved niche did
not outperform specimen-level ecology or a matched shuffle. A variance decomposition supplied the
reason---niche-composition variance, though overwhelmingly cell-local, is stage-invariant at the
available resolution, so the progression-aligned signal it carries resides at the coarser specimen
scale rather than at the individual receiver's neighborhood. A pancreatic precursor series returned a correctly calibrated null,
confirming that the estimator does not manufacture stability where cohort size cannot support it.
StageBridge therefore supports a reproducible, interpretable context-associated component of
premalignant epithelial progression while defining the unresolved spatial scale of that association,
contributing a prespecified estimand with a characterized recovery boundary rather than an unqualified
positive claim.


%% list of keywords seperated by comma
\keywords{premalignant progression, lung adenocarcinoma, spatial transcriptomics, optimal transport,
universal differential equations, gene regulatory networks, tumor microenvironment, single-cell}


%%%%  committee members (add it right after the abstract w/o page break)
\begin{singlespace}

    %% if you have co-advisor, add here w/ \vspace{0.1in} as shown below
    %% alternatively you can use minipage environment to put side-by-side
    \section*{Primary reader and thesis advisor}

    Dr. Chris Bradburne, PhD \\
    Associate Professor\\
    Department of Genetic Medicine\\
    Johns Hopkins University, Baltimore, MD


    \section*{Secondary readers}

    [SECONDARY READER --- NAME TO BE ADDED] \\
    Advanced Academic Programs\\
    Krieger School of Arts and Sciences \\
    Johns Hopkins University, Baltimore, MD

    %% you can add more readers if you have them on your committee 
    %% use \vspace{0.1in} in between members for clarity
    %% you can also place committee members side-by-side using `minipage`


\end{singlespace}